% Generated draft: requires mathematical and source review.
\documentclass[11pt]{article}
\usepackage[margin=25mm]{geometry}
\usepackage{fontspec}
\setmainfont{Latin Modern Roman}
\usepackage{amsmath,amssymb,mathtools,unicode-math}
\setmathfont{Latin Modern Math}
\usepackage{enumitem,needspace,xurl,hyperref}
\setlength{\parindent}{0pt}
\setlength{\parskip}{5pt}
\setlength{\emergencystretch}{3em}
\newcommand{\N}{\mathbb{N}}
\newcommand{\Z}{\mathbb{Z}}
\newcommand{\Q}{\mathbb{Q}}
\newcommand{\R}{\mathbb{R}}
\newcommand{\C}{\mathbb{C}}
\newcommand{\F}{\mathbb{F}}
\newcommand{\A}{\mathbb{A}}
\newcommand{\PP}{\mathbb{P}}
\newcommand{\E}{\mathbb{E}}
\newcommand{\eps}{\varepsilon}
\newcommand{\dd}{\,\mathrm{d}}
\begin{document}
\section*{UnsolvedMath ID: 20000626}
\textbf{Source problem code:} catalogue entry
\subsection{Definitions and mathematical statement}
For a random polynomial \(P_n(z)=\sum_{j=0}^n a_{n,j}z^j\) of degree \(n\geq1\) almost surely, define its largest-root modulus by
\[
R_n=\max\{|z|:P_n(z)=0\}.
\]
For real coefficients this definition still takes all complex roots, not just real roots. Wenbo Li's question asks for asymptotics of \(R_n\) as \(n\to\infty\), for random polynomials in one real or complex variable. The verified source gives no coefficient law, weights, dependence assumptions, or mode of asymptotics. A precise theorem must specify the ensemble and whether it seeks a growth scale, convergence in probability or distribution, tails, or moments. A real-roots-only interpretation would be a different question and would need a convention when no real root exists.

The attached aid illustrates why there is no single ensemble-free scale. For independent standard real or circular complex Gaussian coefficients \(\xi_j\), the Kac ensemble is
\(P_n(z)=\sum_{j=0}^n\xi_jz^j\), and its reported largest-root moduli are tight. The usual Gaussian Weyl ensemble instead has
\(W_n(z)=\sum_{j=0}^n\xi_jz^j/\sqrt{j!}\); the aid reports the lower bound
\[
\Pr\left\{R_n(W_n)\geq(e^{-1/2}-\varepsilon)\sqrt n\right\}\longrightarrow1
\quad\text{for every }\varepsilon>0.
\]
This is a lower bound, not a complete Weyl extreme-root asymptotic.

For a more specific Kac claim, let the real or complex \(\xi_j\) be independent and identically distributed, and suppose that for some \(\alpha>0\), \(u_0,c,C>0\),
\[
c u^\alpha\leq\Pr\{|\xi_0|\leq u\}\leq C u^\alpha
\quad(0<u\leq u_0),\qquad
\mathbb E|\xi_0|^\alpha<\infty.
\]
In particular the coefficient law has no atom at zero. The aid claims constants \(x_0,c_1,C_1>0\), depending on the law but not the degree, such that
\[
c_1x^{-\alpha}\leq\Pr\{R_n>x\}\leq C_1x^{-\alpha}
\quad(n\geq1,\ x\geq x_0).
\]
It reports an associated limiting random-power-series law \(R_\infty\) with the same tail order, and convergence of every subcritical positive moment:
\[
\sup_{n\geq1}\mathbb E R_n^s<\infty,
\qquad \mathbb E R_n^s\longrightarrow\mathbb E R_\infty^s<\infty
\quad(0<s<\alpha).
\]
For \(s\geq\alpha\), the claimed lower tails force divergence of the corresponding moments at every finite degree and for the limiting law. The exact construction of the limiting power-series variable is not reproduced in the record. These are supplied unreviewed, ensemble-specific results; they do not assign an asymptotic to unspecified random polynomials.
\subsection{Short English statement}
Find the size and distribution of the largest-modulus zero as degree grows, after specifying a random-polynomial ensemble. The source specifies none. The supplied aid contrasts tight Gaussian Kac extremes with growing Weyl extremes and gives a conditional Kac tail and moment theorem based on coefficient small-ball behavior.
\subsection{Sources}
\begin{itemize}
\item Wenbo Li, Question 1, AIM workshop \emph{Random analytic functions}: \url{https://aimath.org/WWN/randomzeros/arcc1.pdf}. The supplied verification says the PDF and following questions supply no ensemble definition.
\item Local registry \texttt{data/unsolvedmath/problems.json}, permanent ID \(20000626\), exact source code \texttt{catalogue entry}; canonical \texttt{aim-analysis-notes.json}, notes[157].
\item The attached aid supplies the Kac/Weyl comparison and small-ball theorem without full external bibliography. Attempt 1 is explicitly machine-generated, not peer reviewed, and requires independent verification before citation.
\end{itemize}
\end{document}
